\documentclass[12pt,a4paper]{article}
\usepackage[UTF8]{ctex}
\usepackage{geometry}
\usepackage{graphicx}
\usepackage{booktabs}
\usepackage{amsmath}
\usepackage{listings}
\usepackage{xcolor}
\usepackage{hyperref}
\usepackage{fancyhdr}
\usepackage{titlesec}
\usepackage{enumitem}
\usepackage{tabularx}
\usepackage{float}

\geometry{left=2.5cm, right=2.5cm, top=2.5cm, bottom=2.5cm}

% 代码样式
\lstset{
    language=Java,
    basicstyle=\small\ttfamily,
    keywordstyle=\color{blue},
    commentstyle=\color{gray},
    stringstyle=\color{red},
    numbers=left,
    numberstyle=\tiny\color{gray},
    frame=single,
    breaklines=true
}

% 超链接
\hypersetup{
    colorlinks=true,
    linkcolor=blue,
    urlcolor=blue
}

\title{\textbf{ProcessOS --- 操作系统进程管理模拟器}\\[0.5em]\large 课程设计报告}
\author{操作系统课程设计}
\date{\today}

\begin{document}

\maketitle
\tableofcontents
\newpage

\section{项目概述}

\subsection{项目背景}
操作系统是计算机系统的核心软件，负责管理硬件资源、调度进程运行、处理中断和异常。通过本项目，我们设计并实现了一个完整的操作系统进程管理模拟器，旨在直观展示操作系统的核心工作原理。

\subsection{项目目标}
\begin{enumerate}
    \item 实现一个仿 macOS 风格的 WebOS 桌面系统
    \item 实现 7 种经典进程调度算法的可视化演示
    \item 模拟硬件资源管理（CPU、内存、IO）
    \item 实现进程同步互斥机制
    \item 提供系统事件审计日志
\end{enumerate}

\section{技术架构}

\subsection{技术栈}

\begin{table}[H]
\centering
\begin{tabular}{lll}
\toprule
\textbf{层级} & \textbf{技术} & \textbf{版本} \\
\midrule
前端 & React + Vite + Axios & React 18 / Vite 5 \\
后端 & Java + Spring Boot + JPA & Java 17 / Spring Boot 3.5 \\
数据库 & MySQL + Hibernate & MySQL 8.0 \\
构建工具 & Maven + npm & Maven 3.9 \\
\bottomrule
\end{tabular}
\end{table}

\subsection{系统架构图}
系统采用前后端分离架构，前端通过 RESTful API 与后端通信：

\begin{verbatim}
┌─────────────────────────────────────────────────┐
│           前端 (React + Vite)                     │
│  ┌─────────┐ ┌──────────┐ ┌──────────────────┐  │
│  │ Desktop │ │TaskManager│ │  SchedulerLab    │  │
│  │ (macOS) │ │ (HUD)    │ │  (算法实验室)     │  │
│  └─────────┘ └──────────┘ └──────────────────┘  │
└───────────────────┬─────────────────────────────┘
                    │ REST API
┌───────────────────┴─────────────────────────────┐
│           后端 (Spring Boot)                      │
│  ┌────────────┐ ┌──────────┐ ┌───────────────┐  │
│  │ Controller │ │ Service  │ │  Repository   │  │
│  └─────┬──────┘ └────┬─────┘ └──────┬────────┘  │
│        │             │              │            │
│  ┌─────┴──────┐ ┌────┴─────┐ ┌─────┴─────────┐ │
│  │ Scheduler  │ │ Hardware │ │    MySQL      │ │
│  │ (7种算法)  │ │  Pool   │ │  (持久化)     │ │
│  └────────────┘ └──────────┘ └───────────────┘ │
└─────────────────────────────────────────────────┘
\end{verbatim}

\section{核心功能实现}

\subsection{macOS 桌面系统}

实现了仿 macOS 的桌面环境，包括：

\begin{itemize}
    \item 高清壁纸背景
    \item 顶部菜单栏（半透明毛玻璃效果）
    \item 底部 Dock 栏（悬停放大动效）
    \item 红绿灯窗口控制（左上角关闭/最小化/最大化）
    \item 窗口拖拽、全屏、Dock 自动隐藏
    \item 应用启动器（搜索 + 分类）
\end{itemize}

\subsection{进程调度算法}

实现了 7 种经典调度算法：

\begin{table}[H]
\centering
\begin{tabular}{lll}
\toprule
\textbf{算法} & \textbf{英文名} & \textbf{特点} \\
\midrule
先来先服务 & FCFS & 按到达顺序执行 \\
短作业优先 & SJF & 选执行时间最短的 \\
最短剩余时间优先 & SRTN & SJF 的抢占式版本 \\
时间片轮转 & RR & 每个进程运行固定时间片 \\
非抢占优先级 & Priority & 选优先级最高的 \\
抢占式优先级 & PreemptivePriority & 高优先级可抢占 \\
多级反馈队列 & MFQ & 多级队列 + 时间片递增 \\
\bottomrule
\end{tabular}
\end{table}

\subsection{硬件资源模拟}

\begin{itemize}
    \item \textbf{多核CPU}：8核模拟，支持核心绑定
    \item \textbf{内存管理}：16GB内存池，首次适应算法（First Fit）
    \item \textbf{IO设备}：磁盘读写、网络传输模拟
    \item \textbf{资源抖动}：每秒随机波动，模拟真实负载
\end{itemize}

\subsection{进程同步互斥}

\begin{itemize}
    \item 生产者消费者问题（信号量 + 互斥锁）
    \item 哲学家就餐问题（死锁避免）
    \item 读者写者问题（读写锁）
\end{itemize}

\section{数据库设计}

系统使用 MySQL 数据库持久化以下数据：

\begin{table}[H]
\centering
\begin{tabular}{lll}
\toprule
\textbf{表名} & \textbf{用途} & \textbf{触发时机} \\
\midrule
execution\_log & 进程执行轨迹 & 手动保存/重置时 \\
performance\_metrics & 性能指标统计 & 手动保存/重置时 \\
scenario\_config & 实验场景配置 & 建表时 seed \\
system\_events & 系统事件日志 & 每次操作时 \\
\bottomrule
\end{tabular}
\end{table}

\section{界面设计}

\subsection{桌面系统}
采用 macOS 风格的极简设计：
\begin{itemize}
    \item 深色壁纸 + 半透明菜单栏
    \item 底部 Dock 栏（毛玻璃效果）
    \item 红绿灯窗口控制
    \item 侧滑控制中心
\end{itemize}

\subsection{任务管理器}
采用赛博朋克 HUD 风格：
\begin{itemize}
    \item 圆形仪表盘（CPU/内存使用率）
    \item 波形图（30秒历史趋势）
    \item 核心柱状图（Equalizer）
    \item 内存乐高方块可视化
\end{itemize}

\section{测试与验证}

\subsection{功能测试}

\begin{table}[H]
\centering
\begin{tabular}{lll}
\toprule
\textbf{测试项} & \textbf{预期结果} & \textbf{实际结果} \\
\midrule
启动应用 & 创建进程并分配资源 & 通过 \\
结束进程 & 释放资源并移除进程 & 通过 \\
挂起/恢复 & 正确切换进程状态 & 通过 \\
单实例检测 & 阻止重复启动游戏 & 通过 \\
fork子进程 & 自动创建子进程 & 通过 \\
调度算法切换 & 正确执行不同算法 & 通过 \\
系统日志记录 & 自动记录操作事件 & 通过 \\
\bottomrule
\end{tabular}
\end{table}

\section{总结与展望}

\subsection{项目总结}
本项目成功实现了一个完整的操作系统进程管理模拟器，具有以下特点：
\begin{enumerate}
    \item 仿 macOS 的精美桌面界面
    \item 7 种调度算法的可视化演示
    \item 完整的硬件资源模拟
    \item 实时系统监控和事件日志
\end{enumerate}

\subsection{未来展望}
\begin{enumerate}
    \item 页面置换算法（LRU/FIFO/Clock）
    \item I/O 设备管理（设备申请/阻塞队列）
    \item 死锁检测（银行家算法）
    \item IPC 进程通信（共享内存 + 消息队列）
\end{enumerate}

\end{document}
